\section{Introduction}
\label{sec:intro}

The $\Lambda$CDM concordance model provides an excellent description of the
large-scale matter distribution, yet faces persistent challenges on sub-galactic
scales~\cite{Weinberg2015}.
Among the most prominent is the \emph{core-cusp problem}: $N$-body simulations
of collisionless cold dark matter consistently predict a density profile rising
as $\rho\propto r^{-1}$ towards the centre (a cusp), whereas observations of
dwarf galaxies prefer a nearly constant inner density (a core)~\cite{Oh2015}.
A natural resolution is provided by self-interacting dark matter
(SIDM)~\cite{Spergel2000}, in which DM particles exchange momentum through a
short-range interaction characterised by the cross-section-to-mass ratio
$\sigma/m$.
Heat conduction driven by elastic scattering thermalises the inner halo over a
few Hubble times, converting cusps into isothermal cores whose sizes scale
predictably with halo concentration~\cite{Tulin2018,Kaplinghat2016}.

Classical dwarf-spheroidals (dSphs) orbiting the Milky Way are among the most
dark-matter-dominated systems known, with mass-to-light ratios exceeding
$\Upsilon > 10^{2}\,\msun/L_{\odot}$, and are therefore ideal laboratories for
SIDM phenomenology~\cite{Walker2009}.
Here we use line-of-sight velocity dispersions of eight dSphs from the DESI-DR5
spectroscopic survey to constrain $\sigma/m$ through Jeans-equation modelling.
Throughout we adopt a flat $\Lambda$CDM background with
$H_{0}=67.4\,\mathrm{km\,s^{-1}\,Mpc^{-1}}$ and $\Omega_{m}=0.315$~\cite{Planck2020}.

% ─────────────────────────────────────────────────────────────────────────────
